% =====================================================================
%  Publicación e interoperabilidad — suite de documentos LaTeX
%  Nivel 4 (Avanzado) · clase article.
%  Compilar: latexmk -pdf publicacion.tex
% =====================================================================
\documentclass[11pt,a4paper]{article}

% --- Base estándar ---
\usepackage[utf8]{inputenc}
\usepackage[T1]{fontenc}
\usepackage{lmodern}
\usepackage{microtype}
\usepackage[spanish,es-noshorthands]{babel}
\usepackage{csquotes}
\usepackage{amsmath}\usepackage{amssymb}
\usepackage{booktabs}
\usepackage{tabularx}
\usepackage[margin=2.4cm]{geometry}
\usepackage{enumitem}
\usepackage{xcolor}
\usepackage{listings}
\usepackage{hyperref}
\hypersetup{unicode, pdftitle={Publicación e interoperabilidad en LaTeX},
            pdfauthor={Rodrigo Martín Vidal Galán},
            colorlinks=true, linkcolor=blue!50!black, urlcolor=blue!50!black}
\usepackage[spanish,capitalise,nameinlink]{cleveref}

% --- Estilo de código ---
\definecolor{codeback}{RGB}{245,247,250}
\definecolor{codekw}{RGB}{20,60,120}
\definecolor{codecom}{RGB}{110,120,135}
\lstset{
  basicstyle=\ttfamily\footnotesize,
  commentstyle=\color{codecom}\itshape,
  backgroundcolor=\color{codeback},
  frame=single, rulecolor=\color{codekw!30},
  breaklines=true, showstringspaces=false, columns=fullflexible,
  language=bash, xleftmargin=8pt, framesep=4pt, extendedchars=true,
  literate={á}{{\'a}}1 {é}{{\'e}}1 {í}{{\'i}}1 {ó}{{\'o}}1 {ú}{{\'u}}1
           {ñ}{{\~n}}1 {¿}{{?`}}1 {¡}{{!`}}1,
}
\lstdefinestyle{tex}{language=[LaTeX]TeX, keywordstyle=\color{codekw}\bfseries}
\newcommand{\C}[1]{\texttt{\textbackslash #1}}
\newcommand{\pkg}[1]{\textsf{#1}}

\title{\textbf{Publicación e interoperabilidad}\\[2pt]
       \large Llevar \LaTeX{} al resto del mundo}
\author{Rodrigo Martín Vidal Galán}
\date{Junio de 2026}

\begin{document}
\maketitle

\noindent El documento \LaTeX{} no vive aislado: hay que \textbf{convertirlo} a
otros formatos (un docente pide Word, una web pide HTML), \textbf{prepararlo}
para enviarlo a una revista o repositorio, cuidar la \textbf{calidad y
accesibilidad} del PDF, y a veces \textbf{generar muchos documentos} a partir de
datos. Esta es la capa de \enquote{salida} y puentes hacia afuera.

\newpage
\tableofcontents
\newpage

% =====================================================================
\section{Conversión de formatos: \texttt{pandoc}}
% =====================================================================
La herramienta central de interoperabilidad es \textbf{\texttt{pandoc}}: convierte
entre decenas de formatos (es externo a \LaTeX, se corre en la terminal). Casos
típicos:
\begin{lstlisting}
# LaTeX  ->  Word (.docx)   (lo que mas piden los docentes)
pandoc informe.tex -o informe.docx

# LaTeX  ->  Markdown / HTML
pandoc informe.tex -o informe.md
pandoc informe.tex -o informe.html --mathjax

# Word  ->  LaTeX   (punto de partida desde un .docx ajeno)
pandoc origen.docx -o origen.tex

# usar una bibliografia .bib en la conversion
pandoc informe.tex --citeproc --bibliography=refs.bib -o informe.docx
\end{lstlisting}
\begin{quote}\small
\textbf{Realismo:} pandoc traduce muy bien estructura, listas, tablas y
matemática básica; lo \textbf{muy específico} de \LaTeX{} (TikZ, macros raras,
ciertos paquetes) no siempre sobrevive. Para Word suele requerir un repaso
manual; aun así, ahorra horas.
\end{quote}

% =====================================================================
\section{LaTeX a HTML (matemática en la web)}
% =====================================================================
Para publicar en la web conservando la matemática hay tres caminos:
\begin{center}\small
\renewcommand{\arraystretch}{1.3}
\begin{tabularx}{\linewidth}{@{}lX@{}}
\toprule
\textbf{Herramienta} & \textbf{Notas} \\
\midrule
\texttt{pandoc -{}-mathjax} & Rápido; la matemática la renderiza \textbf{MathJax} en el navegador. \\
\texttt{make4ht} / \texttt{tex4ht} & Más fiel a \LaTeX; genera HTML+CSS (y ODT, ePub). \\
\textbf{LaTeXML} & Conversión de \textbf{alta calidad} a HTML/XML (la usa arXiv para su HTML). \\
\bottomrule
\end{tabularx}
\end{center}
Si solo querés la matemática \emph{en} una web propia, \textbf{MathJax} o
\textbf{KaTeX} renderizan \LaTeX{} directamente en el navegador (pegás
\verb|$...$| en el HTML).

% =====================================================================
\section{Calidad del PDF: metadatos, PDF/A y accesibilidad}
% =====================================================================
El PDF también se \enquote{publica}, y conviene que sea prolijo y duradero.

\paragraph{Metadatos.} Con \pkg{hyperref} (título, autor; aparecen en
\emph{Propiedades} del PDF y ayudan a buscadores/bibliotecas):
\begin{lstlisting}[style=tex]
\hypersetup{unicode, pdftitle={Mi tesina}, pdfauthor={Mi Nombre},
            pdfsubject={...}, pdfkeywords={latex, matematica}}
\end{lstlisting}

\paragraph{PDF/A (archivado).} Muchos repositorios de tesis exigen \textbf{PDF/A}
(formato auto-contenido, con fuentes incrustadas, pensado para preservación a
largo plazo). El paquete \pkg{pdfx} lo activa:
\begin{lstlisting}[style=tex]
\usepackage[a-2b]{pdfx}   % declara conformidad PDF/A-2b
\end{lstlisting}

\paragraph{Accesibilidad (PDF \enquote{etiquetado}).} Un PDF \emph{tagged} permite
que un lector de pantalla lo lea en orden y describa figuras. Es un área en
desarrollo en \LaTeX: el proyecto \enquote{\LaTeX{} Tagged PDF} y el paquete
\pkg{tagpdf} avanzan hacia PDF accesibles; conviene además poner texto
alternativo en las imágenes importantes. Para entregas formales, verificá si te
exigen accesibilidad.

% =====================================================================
\section{Preparar un envío (revista / repositorio / arXiv)}
% =====================================================================
Al enviar a una revista o a \textbf{arXiv} casi siempre piden el \textbf{fuente},
no solo el PDF. Dos cuidados:
\begin{itemize}[leftmargin=*]
  \item \textbf{Aplanar} el proyecto: si usás \C{include}/\C{input}, conviene
        unificar en un solo \texttt{.tex} con \texttt{latexpand}:
\begin{lstlisting}
latexpand tesina.tex > enviar.tex
\end{lstlisting}
  \item \textbf{Bibliografía:} incluí el \texttt{.bbl} generado (arXiv no corre
        biber/bibtex). Y revisá que compile \enquote{en limpio} desde cero.
\end{itemize}
Herramientas como \texttt{arxiv-collector} juntan el \texttt{.tex}, el
\texttt{.bbl} y las figuras en un \texttt{.zip} listo para subir. Cada
revista trae su \textbf{plantilla} (clase + \texttt{.bst}); respetala y leé sus
\emph{author guidelines}.

% =====================================================================
\section{Mail merge: generar muchos documentos desde datos}
% =====================================================================
\LaTeX{} sirve para \textbf{producción en masa}: certificados, exámenes
personalizados, cartas\dots a partir de una tabla de datos. El paquete
\pkg{csvsimple} recorre un CSV y arma un documento por fila:
\begin{lstlisting}[style=tex]
\usepackage{csvsimple}
...
\csvreader[head to column names]{alumnos.csv}{}{%
  \newpage
  Certificamos que \nombre\ \apellido\ aprobo con \nota.%
}
\end{lstlisting}
Con un CSV de 100 filas, una sola compilación genera las 100 páginas. La
alternativa más potente es \pkg{datatool} (ordenar, filtrar, hacer cuentas sobre
los datos). Combinado con \pkg{pgfplots}, también graficás esos datos.

% =====================================================================
\section{Compartir y colaborar}
% =====================================================================
\begin{itemize}[noitemsep,leftmargin=*]
  \item \textbf{Overleaf:} edición en la nube, sin instalar, con historial y
        comentarios; ideal para co-autoría y para enviar a quien no tiene \LaTeX.
  \item \textbf{git:} el fuente es texto, versionalo (ver \emph{Flujo}); para que
        un docente vea cambios, \texttt{latexdiff} marca lo agregado/borrado.
  \item \textbf{El PDF} es el entregable universal: se ve igual en todos lados.
        Cuando piden \enquote{editable}, ahí entra pandoc $\to$ Word.
\end{itemize}

% =====================================================================
\section{Buenas prácticas}
% =====================================================================
\begin{itemize}[noitemsep,leftmargin=*]
  \item Si te piden Word/HTML, \texttt{pandoc} primero y \textbf{repasá} el
        resultado; no esperes una conversión perfecta de lo muy específico.
  \item Para tesis/repositorios, averiguá si exigen \textbf{PDF/A} (y usá
        \pkg{pdfx}) y/o accesibilidad.
  \item Antes de enviar: \textbf{aplaná} (\texttt{latexpand}), incluí el
        \texttt{.bbl}, y verificá una compilación limpia desde cero.
  \item Metadatos del PDF siempre (\pkg{hyperref}); \texttt{unicode} activado.
  \item Producción en masa $\to$ \pkg{csvsimple}/\pkg{datatool} desde un CSV.
\end{itemize}

% =====================================================================
\section*{Ver también}
% =====================================================================
\addcontentsline{toc}{section}{Ver también}
\begin{itemize}[noitemsep,leftmargin=*]
  \item \emph{Flujo e instalación} — git, \texttt{latexdiff}, Overleaf, CI.
  \item \emph{Referencias y bibliografía} — el \texttt{.bbl} y los estilos.
  \item \emph{Programación} — generar \texttt{.tex} desde datos/otros lenguajes.
  \item Externo: \texttt{pandoc.org}; \texttt{texdoc pdfx}, \texttt{texdoc csvsimple}.
\end{itemize}

\end{document}
